% One change at a time: the runs that change one thing in another run, set
% against that run, and how to read them. The same for every model, so a
% paper reads this under its own heading, after its results, then says what
% its own numbers mean, then reads studies.tex. Each part waits for its runs:
% a model that has none prints nothing.
\ifpairs
\paragraph{Runs set against the run each changed.}
The ladder says what each learner is worth. A change asks a narrower
question: what is one thing worth inside a learner? Twice the epochs, a body
that starts from another run's weights, the hand features under a network.
Each is a pair of runs that differ in that one thing.
Table~\ref{tab:pairs} sets the first of each pair against the second on the
same resampled test games, the way every run is set against the reference,
and gives the first run's gain. This is where the lab's rule for keeping a
change is read: a change is kept as a run of the model's own when its run is
better on the valid rows and the interval of its gain on the test rows
leaves zero out, however good its own score looks. A change that is not kept
is not thrown away: its row stays in this table, as a measurement of what
the change is worth.

\begin{table}[htbp]\centering\footnotesize
\block{pairs}
\caption{Runs set against the run each changed or has to beat. A gain is
the first run's over the second's, positive when the first is better. On the
valid rows it is the difference of the two scores. On the test rows it is
measured on the same resampled games, and it is bold when its 95\% interval
leaves zero out.}
\label{tab:pairs}
\end{table}
\fi

\ifvariants
Figure~\ref{fig:variants} follows each network of the ladder that another
run changed, beside the runs that changed it, on the valid rows. A run keeps
the colour of the network it changed, and its line is broken by what it
changed: long dashes for a longer schedule, dots for a start from another
run's body.

\begin{figure}[htbp]\centering
\block{variants}
\caption{A network of the ladder (the unbroken line) and the runs that
changed one thing in it (the broken lines), after every epoch on the valid
rows. One panel a network, all on one scale. The dashed level is the
reference's score on the same rows.}
\label{fig:variants}
\end{figure}
\fi
